\subsection{Conclusions of the Sensitivity Analysis}
\label{sec:p3-conclusions}
The separation inventory is conditional on the selected propulsion, aerodynamics, recovery plan, and retained propellant. For the reference design, the P2 candidate carries approximately 26.23 t, including the 5 t planned usable reserve, with approximately 5.83 t methane and 20.40 t oxygen. The known-efficiency designs give accepted local inventories of 24.20--28.78 t. The higher-crossflow design requires approximately 27.13 t, while the lower-crossflow search remains unresolved. These individual designs use separately selected ship positions and do not define a common inventory or one command history valid across uncertain conditions.

At a fixed ship, the original nominal commands are sensitive to initial and wind errors. Case-specific separation-time planning restores the eleven examined cases at a common 30 t trial inventory when the initial state and steady wind are accurately known. With wind unknown to the planner, the 20 ms feedback scheme restores the five examined points from $-0.1$ to $+0.1$ m/s and retains more than 6.03 t. This sampled result supports local correction under exact state information and ideal actuator response; it does not establish a continuous wind envelope or robustness to combined errors.

Touchdown recovery and low flight loads are separate requirements. The nominal feedback trajectory reaches approximately 304 kPa and $8.13\,g_0$, above the corresponding P2 peaks. Pre-ignition pressure and flow-energy peaks cannot be reduced by changing only the later feedback. Structural and thermal safety remain undetermined without vehicle allowables and material-response models. The present results therefore provide conditional propellant budgets and specific limits of the recovery calculations. Additional thrust, reserve, deck-motion, actuator, and three-dimensional studies are possible extensions, rather than prerequisites for the bounded sensitivity analysis reported here.
